% ============================================================
%  Black-Litterman with Ledoit-Wolf Shrinkage
%  Fixed-Omega leakage and the critical correlation threshold
%  Standalone paper draft - Lucas Posern
%
%  All quoted numbers are reproduced by code/paper_numbers.py from
%  the frozen data in ../data/ (thesis pipeline: excess returns,
%  sample cov ddof=1, sklearn Ledoit-Wolf, implied delta).
%  Revision 2026-07: corrected rho* (factor-2 bug), exact per-pair
%  Delta c_eff table, updated backtest to thesis view set, paper
%  frame (abstract, introduction, conclusion), decimal points.
% ============================================================

\documentclass[a4paper,11pt]{article}

\usepackage[a4paper, left=2.5cm, right=2.5cm, top=2.5cm, bottom=2.5cm]{geometry}
\usepackage{amsmath, amssymb, amsthm, amsfonts}
\usepackage{mathtools}
\usepackage[english]{babel}
\usepackage[T1]{fontenc}
\usepackage[utf8]{inputenc}
\usepackage{booktabs}
\usepackage{array}
\usepackage{float}
\usepackage{graphicx}
\usepackage{caption}
\usepackage{xcolor}
\usepackage[round]{natbib}
\usepackage{listings}
\usepackage[colorlinks=true, linkcolor=blue!70!black, citecolor=blue!70!black, urlcolor=blue!70!black]{hyperref}

\lstset{
  language=Python,
  basicstyle=\ttfamily\footnotesize,
  keywordstyle=\color{blue!70!black}\bfseries,
  commentstyle=\color{green!40!black}\itshape,
  stringstyle=\color{red!70!black},
  breaklines=true,
  columns=fullflexible,
  showstringspaces=false,
  frame=single,
  framerule=0.4pt,
  rulecolor=\color{gray!60},
  xleftmargin=4pt,
  xrightmargin=4pt,
}

% ------------------------------------------------------------
% Theorem environments
% ------------------------------------------------------------
\theoremstyle{plain}
\newtheorem{proposition}{Proposition}
\newtheorem{theorem}[proposition]{Theorem}
\newtheorem{corollary}[proposition]{Corollary}
\theoremstyle{definition}
\newtheorem{example}[proposition]{Example}
\newtheorem{remark}[proposition]{Remark}

% ------------------------------------------------------------
% Notation
% ------------------------------------------------------------
\newcommand{\R}{\mathbb{R}}
\newcommand{\E}{\mathbb{E}}
\newcommand{\Sig}{\Sigma}
\newcommand{\Om}{\Omega}
\newcommand{\Mprec}{M}
\newcommand{\muBL}{\mu_{\mathrm{BL}}}
\newcommand{\Std}{\mathrm{Std}}
\newcommand{\LW}{\mathrm{LW}}
\newcommand{\eff}{\mathrm{eff}}

\title{Black--Litterman with Ledoit--Wolf Shrinkage:\\
Fixed-$\Omega$ Leakage and the Critical Correlation Threshold%
\thanks{This paper grew out of the author's bachelor thesis at TU Dresden.
I thank Rodrigo S.~Targino (FGV EMAp, Rio de Janeiro) for detailed feedback
on the proof strategy. All remaining errors are mine.}}
\author{Lucas Posern\\
\small Faculty of Mathematics, TU Dresden\\
\small \texttt{posernlucas@gmail.com}}
\date{\today}

% ------------------------------------------------------------
\begin{document}

\maketitle

\begin{abstract}
\noindent
The Black--Litterman model is routinely combined with shrinkage estimators
of the covariance matrix, most prominently the Ledoit--Wolf estimator.
We show that the effect of the estimator switch on the posterior depends
entirely on how the view-uncertainty matrix $\Omega$ is handled. Under
Idzorek's confidence calibration, $\Omega$ co-moves with $\Sigma$ and the
switch is exactly inert in the view channel. If instead $\Omega$ is
calibrated once on the sample covariance and then frozen---a workflow that
appears regularly in practice---the effective view confidence
$c_{\eff}$ shifts by a closed-form amount that is independent of the prior
scale $\tau$ and grows with the shrinkage intensity. The shift is positive
precisely when the view variance lies below twice the average asset
variance, a condition met by relative views on highly correlated pairs
\emph{and} on low-variance (bond) pairs. We derive an exact expression for
the induced changes in posterior means and weights under a single relative
view, provide a closed-form critical correlation threshold, and quantify
the effect on the eight-asset Idzorek ETF universe: at a shrinkage
intensity of $0.041$, five of 28 pairs shift their effective confidence by
more than five percentage points, the largest by $+10.6$. Recomputing
$\Omega$ alongside $\Sigma$ removes the leakage exactly.
\end{abstract}

\medskip
\noindent\textbf{Keywords:} Black--Litterman; Ledoit--Wolf shrinkage;
view uncertainty; covariance estimation; portfolio construction.\\
\textbf{JEL:} G11; C13; C58.

% ============================================================
\section{Introduction}
\label{sec:intro}
% ============================================================

The Black--Litterman model \citep{BlackLitterman1992} is one of the few
pieces of portfolio theory that made the crossing from journals to
practice. It anchors expected returns at the market equilibrium and blends
in subjective views with an explicit uncertainty matrix $\Omega$, thereby
taming the input sensitivity of mean--variance optimisation documented by
\citet{Michaud1989} and \citet{BestGrauer1991}. Its second input, the
covariance matrix $\Sigma$, receives far less attention in the
Black--Litterman literature, although the sample estimator is unreliable
whenever the number of assets is not small relative to the sample size.
The standard remedy is the shrinkage estimator of \citet{LedoitWolf2004},
which pulls the sample covariance toward a scaled identity and improves
the conditioning of $\Sigma^{-1}$ by orders of magnitude in equity
universes \citep{LedoitWolfHoney2004}.

This paper studies what happens inside the Black--Litterman model when the
covariance estimator is switched from the sample estimator to Ledoit--Wolf.
The answer depends entirely on a bookkeeping question that, to our
knowledge, has not been isolated in the literature: what happens to
$\Omega$? Two workflows coexist in practice. In the first, $\Omega$ is
derived from $\Sigma$ through Idzorek's confidence calibration
\citep{Idzorek2007} and therefore co-moves with the estimator. In the
second, $\omega_k$ is computed once on the sample covariance, stored as a
set of scalars, and left untouched when the covariance model is later
upgraded---for instance because the badly conditioned
$\hat\Sigma_s^{-1}$ produced visibly implausible leverage.

Our contribution is threefold. First, we show that under Idzorek's
calibration the estimator switch is \emph{exactly} inert in the view
channel: the effective view absorption equals the specified confidence
$c$ for any positive-definite estimator, and the entire effect on the
posterior runs through the equilibrium prior, bounded at
$O(\alpha)$. Second, for the fixed-$\Omega$ workflow we derive a
closed-form expression for the shift $\Delta c_{\eff}$ of the effective
view confidence. The expression is exact, independent of $\tau$, linear in
the shrinkage intensity $\alpha$ for small $\alpha$, and its sign is
determined by a simple criterion: the view variance versus twice the
average asset variance. In the symmetric case of equal asset volatilities
this criterion translates into a closed-form critical correlation
threshold $\rho^*(\alpha, \varepsilon)$. An exact identity for the
Black--Litterman posterior under a single relative view then propagates
$\Delta c_{\eff}$ analytically into posterior means and portfolio
weights. Third, we quantify the effect on the eight-asset ETF universe of
\citet{Idzorek2007}: with the analytic shrinkage intensity
$\alpha = 0.041$, five of the 28 asset pairs shift their effective
confidence by more than five percentage points---among them a bond pair
with moderate correlation, which shows that the popular intuition
``only high correlations are affected'' is incomplete. An out-of-sample
backtest separates this view channel from the precision channel through
$\Sigma^{-1}$.

The practical lesson is procedural rather than statistical. Covariance
estimator, confidence and view uncertainty form one unit; upgrading one
piece without re-deriving the others silently moves the model away from
the analyst's stated inputs. Related observations---that constraints or
estimators interact with estimation error in ways invisible at the input
stage---appear in \citet{JagannathanMa2003} for long-only constraints and
in \citet{DeMiguelGarlappiUppal2009} for the difficulty of beating naive
diversification out of sample. Within the Black--Litterman literature,
\citet{HeLitterman1999}, \citet{SatchellScowcroft2000} and
\citet{Walters2014} document the model's sensitivity to the
$(\tau, \Omega)$ specification; our result gives one mechanism of this
sensitivity a closed form. On the shrinkage side, \citet{LedoitWolf2003}
and the nonlinear extension of \citet{LedoitWolf2017} would enter our
formulas only through a different expression for the view-variance shift;
the leakage mechanism itself is target-independent.

The paper proceeds as follows. Section~\ref{sec:blm-brief} fixes notation,
model and data. Sections~\ref{sec:sample-problem} and~\ref{sec:lw} review
the conditioning problem and the Ledoit--Wolf estimator.
Section~\ref{sec:idzorek-cancel} settles the Idzorek case.
Section~\ref{sec:fixed-omega} contains the main results for fixed
$\Omega$. Section~\ref{sec:precision} isolates the precision channel
through $\Sigma^{-1}$, Section~\ref{sec:backtest} reports the
out-of-sample backtest, and Section~\ref{sec:conclusion} concludes.

% ============================================================
\section{The Black--Litterman Model and the Data}
\label{sec:blm-brief}
% ============================================================

For $n$ risky assets with covariance $\Sigma\in\R^{n\times n}$,
$\Sigma\succ 0$, the Black--Litterman posterior mean is
\begin{equation}
\label{eq:bl-posterior}
  \muBL \;=\; \bigl[(\tau\Sig)^{-1} + P^\top\Om^{-1}P\bigr]^{-1}
              \bigl[(\tau\Sig)^{-1}\Pi \;+\; P^\top\Om^{-1}Q\bigr],
\end{equation}
with implicit equilibrium return $\Pi = \delta\Sigma w_{\mathrm{eq}}$,
view-pick matrix $P\in\R^{K\times n}$, view returns $Q\in\R^K$,
view-uncertainty matrix $\Om\succ 0$, prior scale $\tau>0$, and risk
aversion $\delta>0$. The unconstrained portfolio is
$w_{\mathrm{BL}} = \tfrac{1}{\delta}\Sigma^{-1}\muBL$.

\paragraph{Data.}
We use the eight-asset ETF universe of \citet{Idzorek2007}: AGG (US
bonds), BWX (international bonds), IWF/IWD (US large-cap growth/value),
IWO/IWN (US small-cap growth/value), EFA (international developed
equity) and EEM (emerging markets equity). Monthly total returns are
taken from Yahoo Finance (adjusted closes) and converted to excess
returns with the 13-week T-bill rate; the series are complete over the
whole analysis period. The training window is January 2008 to December
2018 ($T=132$, $n=8$); January 2019 to December 2024 is reserved for the
out-of-sample test in Section~\ref{sec:backtest}. The market weights
$w_{\mathrm{eq}}$ are the market-capitalisation weights published by
Idzorek for this universe and enter only as the prior anchor. All data,
the estimation pipeline and every number quoted in this paper are
reproduced by a single script; see the reproducibility statement at the
end.

\paragraph{Calibrated quantities.}
On the training window the implied risk aversion of the S\&P~500 excess
return series is $\delta = 2.44$; we use $\tau = 0.05$ and view
confidences $c_k = 0.50$ throughout. The average asset variance, which
acts as the Ledoit--Wolf shrinkage target scale, is
$\mu_s = \tfrac{1}{n}\mathrm{tr}(\hat\Sigma_s) = 0.0282$ annualised
($2.35\cdot 10^{-3}$ per month), and the analytic shrinkage intensity is
$\alpha = 0.041$. Three relative views with these confidences run through
the multi-view sections: IWF $>$ IWD by $+3\,\%$ p.a., EEM $>$ EFA by
$+4\,\%$ p.a., and IWF $>$ AGG by $+6\,\%$ p.a. The single-view analysis
of Section~\ref{sec:fixed-omega} uses the first of these.

% ============================================================
\section{The Sample Covariance Problem}
\label{sec:sample-problem}
% ============================================================

The sample covariance estimator is
\[
  \hat\Sigma_s \;=\; \frac{1}{T-1}\sum_{t=1}^T (r_t - \bar r)(r_t - \bar r)^\top.
\]
For $n$ assets and $T$ observations this is reliable only if $T \gg n$.
Once $n/T$ is not negligible, the eigenvalues of $\hat\Sigma_s$ spread
too widely; large eigenvalues are overestimated, small ones are
underestimated \citep[p.~371]{LedoitWolf2004}. The smallest eigenvalue
moves toward zero, the condition number
$\kappa(\Sigma) = \lambda_{\max}/\lambda_{\min}$ grows, and the inverse
$\hat\Sigma_s^{-1}$ used in
$w_{\mathrm{BL}} = \tfrac{1}{\delta}\Sigma^{-1}\muBL$ amplifies
estimation error. On the training window the condition number is
$\kappa(\hat\Sigma_s) = 341.8$.

% ============================================================
\section{Ledoit--Wolf Shrinkage}
\label{sec:lw}
% ============================================================

\citet{LedoitWolf2004} regularise $\hat\Sigma_s$ by shrinking it toward
the scaled identity $\mu_s I_n$, with
$\mu_s = \tfrac{1}{n}\mathrm{tr}(\hat\Sigma_s)$,
\[
  \hat\Sigma_{\LW} \;=\; (1-\alpha)\,\hat\Sigma_s \;+\; \alpha\,\mu_s\,I_n,
\]
where the shrinkage coefficient $\alpha\in[0,1]$ minimises the expected
squared Frobenius error $\E[\|\hat\Sigma - \Sigma\|_F^2]$. Because
$\alpha\mu_s I_n$ adds at least $\alpha\mu_s>0$ to every eigenvalue,
$\hat\Sigma_{\LW}\succ 0$ is automatic. On the training window the
analytic estimator yields $\alpha = 0.041$ and brings the condition
number down to $\kappa(\hat\Sigma_{\LW}) = 108.0$, a factor $3.2$
reduction.

% ============================================================
\section{Idzorek $\Omega$: the Estimator Switch is Inert}
\label{sec:idzorek-cancel}
% ============================================================

We can keep this case brief. When $\Omega$ is derived from $\Sigma$ via
Idzorek's formula, the estimator cancels exactly in the view channel.

Idzorek sets
\[
  \omega_{kk} \;=\; \frac{1-c_k}{c_k}\,\tau\,p_k^\top\Sigma p_k,
\]
and the posterior simplifies along the view direction to
\[
  P\muBL \;=\; (1 - c)\,P\pi \;+\; c\,Q.
\]
The view absorption equals $c$ exactly, independent of $\Sigma$.
Switching $\hat\Sigma_s \to \hat\Sigma_{\LW}$ scales the view variance
and the matching $\omega$ proportionally, the ratio stays at $c$. The
entire LW impact on the posterior runs through the prior, and is bounded
by $\Delta\pi = O(\alpha)$.

Concretely, the LW estimator shifts the prior linearly,
\[
  \Delta\pi_i \;=\; \alpha\,\delta\,\Bigl[(\mu_s - \sigma_i^2)\,w_{\mathrm{eq},i}
        \;-\; \sum_{j\neq i}\hat\Sigma_{s,ij}\,w_{\mathrm{eq},j}\Bigr].
\]
Assets with above-average variance ($\sigma_i^2 > \mu_s$) receive a lower
prior, assets with below-average variance receive a higher one. On the
Idzorek universe $|\Delta\pi_i| \leq 0.25\,\%$ p.a.\ for all eight ETFs.

The Idzorek world is settled. The interesting case follows.

% ============================================================
\section{Fixed $\Omega$: What Happens?}
\label{sec:fixed-omega}
% ============================================================

A different workflow appears regularly in practice. The analyst computes
$\omega_k$ once on the sample estimator and stores the scalars. Later,
the covariance is replaced by an LW estimator (perhaps because the badly
conditioned $\hat\Sigma_s^{-1}$ produced obviously implausible leverage),
but $\omega_k$ stays fixed. The view absorption then shifts implicitly,
because $\Sigma$ moves in the numerator while $\omega_k$ does not move
in the denominator.

We derive the closed form for the shift in $c_{\eff}$, translate it into
a general sign-and-size criterion and, for the symmetric case, into a
critical correlation threshold, and verify the chain on the IWF/IWD
pair. Throughout this section the analysis rests on a single relative
view; the symmetric per-component statements additionally assume
$\sigma_i \approx \sigma_j$ for the view pair, an assumption we make
explicit where it is used and drop where it is not needed.

\subsection{Effective View Absorption $c_{\eff}$}
\label{sec:ceff}

The actual view absorption after evaluating the posterior is
\[
  c_{\eff,k} \;=\; \frac{\tau\,v_k}{\tau\,v_k + \omega_k},
  \qquad v_k := p_k^\top\Sigma p_k.
\]
$c_{\eff,k}$ measures how far the posterior travels along the view
direction from prior toward view. At $c_{\eff,k} = 0$ the posterior
stays at the prior, at $c_{\eff,k} = 1$ it lands exactly on the view.
The ratio depends on two quantities, the view variance $v_k$ and the
view uncertainty $\omega_k$. The role of $\tau$ matters here. $\tau$
controls how tightly the prior is held; small $\tau$ pulls $c_{\eff,k}$
down, large $\tau$ lets the posterior move toward the view more easily.

In the Idzorek world with $\omega_k = \tfrac{1-c_k}{c_k}\tau v_k$, both
$\tau$ and $v_k$ cancel and $c_{\eff,k} = c_k$ exactly. In the
fixed-$\Omega$ world this cancellation breaks.

\subsection{Shift of $v_k$ Under Shrinkage}
\label{sec:vk-shift}

For a relative view $p_k = e_i - e_j$ the view variance is
$v_k = \sigma_i^2 + \sigma_j^2 - 2\hat\Sigma_{s,ij}
     = \sigma_i^2 + \sigma_j^2 - 2\rho_{ij}\sigma_i\sigma_j$,
with no assumption on the volatilities. Under LW,
\[
  v_k^{\LW} \;=\; (1 - \alpha)\,v_k^\Std \;+\; 2\alpha\,\mu_s,
\]
since $p_k^\top I_n p_k = 2$. Hence $v_k$ grows under shrinkage exactly
when $v_k^\Std < 2\mu_s$, i.e.\ when the view variance lies below twice
the average asset variance. For positively correlated pairs of similar
volatility this is the typical case; it also holds for pairs whose
variances are simply small relative to the universe average, such as
bond pairs. With $v_k$, the absorption $c_{\eff,k}$ grows.

\subsection{Closed Form for $\Delta c_{\eff,k}$}
\label{sec:dceff}

Combining $c_{\eff,k}(\hat\Sigma_{\LW})$ and $c_{\eff,k}(\hat\Sigma_s)$
over a common denominator and simplifying,
\begin{align*}
  \Delta c_{\eff,k}
  \;&=\; \frac{\tau\,v_k^\LW}{\tau\,v_k^\LW + \omega_k}
        \;-\; \frac{\tau\,v_k^\Std}{\tau\,v_k^\Std + \omega_k}\\
  \;&=\; \frac{\tau\,v_k^\LW\,(\tau\,v_k^\Std + \omega_k) \;-\; \tau\,v_k^\Std\,(\tau\,v_k^\LW + \omega_k)}
              {(\tau\,v_k^\LW + \omega_k)(\tau\,v_k^\Std + \omega_k)}\\
  \;&=\; \frac{\tau\,\omega_k\,(v_k^\LW - v_k^\Std)}
              {(\tau\,v_k^\LW + \omega_k)(\tau\,v_k^\Std + \omega_k)}.
\end{align*}
Substituting $v_k^\LW - v_k^\Std = \alpha(2\mu_s - v_k^\Std)$ yields
\begin{equation}
\label{eq:dceff}
  \Delta c_{\eff,k}
  \;=\;
  \frac{\tau\,\omega_k\,\alpha\,(2\mu_s - v_k^\Std)}
       {(\tau\,v_k^\LW + \omega_k)(\tau\,v_k^\Std + \omega_k)}.
\end{equation}
The shift scales linearly in $\alpha$ and $\omega_k$, and its sign
matches $\mathrm{sgn}(2\mu_s - v_k^\Std)$: the general criterion
announced above, valid for arbitrary volatility ratios.

In the workflow under study, $\omega_k$ was calibrated on the sample
estimator, $\omega_k = \tfrac{1-c}{c}\,\tau\,v_k^\Std$. Substituting
this into \eqref{eq:dceff} cancels $\tau$ completely and gives the exact
form
\begin{equation}
\label{eq:dceff-exact}
  \Delta c_{\eff,k}
  \;=\;
  \frac{c\,(1-c)\,\bigl(v_k^\LW - v_k^\Std\bigr)}
       {c\,v_k^\LW + (1-c)\,v_k^\Std}.
\end{equation}
Two consequences are worth recording. First, the leakage does not depend
on the prior scale $\tau$ at all; no choice of $\tau$ mitigates it.
Second, \eqref{eq:dceff-exact} is exact for any pair of positive view
variances, so it can be evaluated asset pair by asset pair without any
symmetry assumption. As $v_k^\Std \to 0$ at fixed $\mu_s$ the shift
approaches its supremum $1-c$.

\subsection{Critical Correlation $\rho^*$ in the Symmetric Case}
\label{sec:rho-star}

To turn \eqref{eq:dceff-exact} into a rule of thumb, consider the
symmetric case $\sigma_i = \sigma_j = \sigma$, in which
$v_k^\Std = 2\sigma^2(1 - \rho_{ij})$. Setting
$\Delta c_{\eff,k} = \varepsilon$ in \eqref{eq:dceff-exact} and solving
for $\rho_{ij}$ with $v_k^\LW = (1-\alpha)v_k^\Std + 2\alpha\mu_s$
yields the closed-form threshold
\begin{equation}
\label{eq:rho-star}
  \rho^*(\alpha,\varepsilon)
  \;=\;
  1 \;-\; \frac{c\,\alpha\,\mu_s\,[(1-c) - \varepsilon]}
              {\sigma^2\,[\varepsilon\,(1 - c\,\alpha) \;+\; c\,(1-c)\,\alpha]},
  \qquad \varepsilon \in (0,\,1-c).
\end{equation}
A view triggers $|\Delta c_{\eff,k}| > \varepsilon$ exactly when
$\rho_{ij} > \rho^*$. The threshold falls as $\alpha$ rises and as the
tolerated shift $\varepsilon$ falls. The symmetric case is a structural
sketch rather than the general statement: for $\sigma_i \neq \sigma_j$
the exact criterion remains \eqref{eq:dceff-exact} evaluated at the
heterogeneous $v_k^\Std = \sigma_i^2 + \sigma_j^2 -
2\rho_{ij}\sigma_i\sigma_j$, and we use that exact form for the
pair-level results below.

\subsection{Values and Economic Statement}
\label{sec:values}

Table~\ref{tab:rho-star} shows $\rho^*$ for four shrinkage intensities
and four significance thresholds. The parameters $\mu_s = 4.0\cdot
10^{-3}$, $\sigma^2 = 3.5\cdot 10^{-3}$, $c = 0.50$ are illustrative,
chosen of the same order as the monthly quantities in our data
($\mathrm{tr}(\hat\Sigma_s)/n = 2.35\cdot 10^{-3}$ per month); the table
depends on the parameters only through the ratio $\mu_s/\sigma^2$
and $c$.

\begin{table}[H]
\centering
\caption{Critical correlation $\rho^*(\alpha,\varepsilon)$ from
\eqref{eq:rho-star}. The pair correlation $\rho_{ij}$ must exceed this
value for $|\Delta c_{\eff}| > \varepsilon$. Illustrative parameters
$\mu_s = 4.0\cdot 10^{-3}$, $\sigma^2 = 3.5\cdot 10^{-3}$, $c = 0.50$.}
\label{tab:rho-star}
\begin{tabular}{lcccc}
\toprule
$\varepsilon$ & $\alpha = 0.04$ & $\alpha = 0.10$ & $\alpha = 0.166$ & $\alpha = 0.30$ \\
\midrule
$1\,\%$  & $0.434$ & $0.188$ & $0.083$ & $-0.006$ \\
$2\,\%$  & $0.629$ & $0.377$ & $0.239$ & $0.106$ \\
$5\,\%$  & $0.826$ & $0.645$ & $0.511$ & $0.343$ \\
$10\,\%$ & $0.915$ & $0.810$ & $0.715$ & $0.571$ \\
\bottomrule
\end{tabular}
\end{table}

At moderate shrinkage the threshold is demanding only for tight
tolerances: with $\alpha = 0.04$, a five-percentage-point shift already
occurs above $\rho^* = 0.826$, a correlation level routine for style
pairs within one equity market. At strong shrinkage ($\alpha = 0.30$),
every non-negatively correlated pair crosses the one-percent band. The
intensity $\alpha$ itself grows with $n/T$, so larger universes on the
same window sit further to the right in the table.

The symmetric threshold understates the reach of the effect, because the
general criterion is $v_k^\Std < 2\mu_s$, not high correlation per se.
Table~\ref{tab:pairs} evaluates the exact expression
\eqref{eq:dceff-exact} for every pair of the universe, with the full
$8\times 8$ covariance shrunk once at $\alpha = 0.041$ (the
practitioner workflow) and $c = 0.50$. Five of 28 pairs shift by more
than five percentage points. The top of the list is dominated by the
highly correlated, nearly homogeneous style pairs, led by IWF/IWD at
$+10.6$ percentage points---the model operates at an effective
confidence of $0.61$ where the analyst specified $0.50$. But the bond
pair AGG/BWX ranks third at $+8.1$ percentage points despite a moderate
correlation of $0.57$: its view variance is small relative to $2\mu_s$
because bond variances are far below the universe average. A threshold
stated in correlations alone would miss it.

\begin{table}[H]
\centering
\caption{Exact $\Delta c_{\eff}$ from \eqref{eq:dceff-exact} for all
pairs with a shift above five percentage points. Full-universe shrinkage
at $\alpha = 0.041$, $\omega$ frozen on the sample estimator,
$c = 0.50$; $k = \sigma_i/\sigma_j$. Training window 2008--2018.}
\label{tab:pairs}
\begin{tabular}{lccc}
\toprule
Pair & $\rho_{ij}$ & $k$ & $\Delta c_{\eff}$ \\
\midrule
IWF/IWD & $0.920$ & $0.99$ & $+10.6$ pp \\
IWO/IWN & $0.939$ & $1.03$ & $+8.7$ pp \\
AGG/BWX & $0.574$ & $0.45$ & $+8.1$ pp \\
IWD/IWN & $0.917$ & $0.80$ & $+6.4$ pp \\
IWF/IWO & $0.909$ & $0.77$ & $+5.4$ pp \\
\bottomrule
\end{tabular}
\end{table}

Translated into economic language: whenever a relative view connects two
assets whose spread variance is small---because they are highly
correlated, or because both carry little variance to begin with---a
frozen $\Omega$ lets the shrinkage-induced widening of the spread
variance flow one-for-one into extra view confidence. The model believes
the view more strongly than the analyst stated. For weakly correlated,
high-variance pairs or negligible $\alpha$, $c_{\eff}$ stays essentially
put and the estimator switch is harmless.

\subsection{From $\Delta c_{\eff}$ to $\Delta\mu$ and $\Delta w$}
\label{sec:dmu-dw}

The shifts $\Delta\mu$ and $\Delta w$ follow analytically from
$\Delta c_{\eff}$. The key step is an exact closed form for the
posterior under a single relative view, which we derive directly from
the posterior equation.

With a single view, $P = p^\top$ and $\Omega = \omega$ is a scalar; the
BL posterior equation reads
\begin{align*}
  M\,\muBL \;=\; (\tau\Sigma)^{-1}\pi \;+\; \tfrac{1}{\omega}\,p\,Q,
  \qquad M \;=\; (\tau\Sigma)^{-1} \;+\; \tfrac{1}{\omega}\,p\,p^\top.
\end{align*}
Left-multiplication by $\tau\Sigma$ and solving for $\muBL$:
\begin{align*}
  \tau\Sigma\,M\,\muBL \;&=\; \tau\Sigma\,\bigl[(\tau\Sigma)^{-1}\pi + \tfrac{1}{\omega}\,p\,Q\bigr]\\
  \Bigl[I \;+\; \tfrac{\tau}{\omega}\,\Sigma p\,p^\top\Bigr]\,\muBL \;&=\; \pi \;+\; \tfrac{\tau Q}{\omega}\,\Sigma p\\
  \muBL \;+\; \tfrac{\tau\,(p^\top\muBL)}{\omega}\,\Sigma p \;&=\; \pi \;+\; \tfrac{\tau Q}{\omega}\,\Sigma p\\
  \muBL \;&=\; \pi \;+\; \tfrac{\tau\,(Q - p^\top\muBL)}{\omega}\,\Sigma p.
  \tag{$\star$}\label{eq:bl-implicit}
\end{align*}

Applying $p^\top$ to $(\star)$ produces a scalar equation in
$p^\top\muBL$ with $v_k := p^\top\Sigma p$:
\begin{align*}
  p^\top\muBL \;&=\; p^\top\pi \;+\; \tfrac{\tau v_k\,(Q - p^\top\muBL)}{\omega}\\
  \omega\,p^\top\muBL \;&=\; \omega\,p^\top\pi \;+\; \tau v_k\,Q \;-\; \tau v_k\,p^\top\muBL\\
  (\omega + \tau v_k)\,p^\top\muBL \;&=\; \omega\,p^\top\pi \;+\; \tau v_k\,Q\\
  p^\top\muBL \;&=\; \tfrac{\omega}{\omega + \tau v_k}\,p^\top\pi \;+\; \tfrac{\tau v_k}{\omega + \tau v_k}\,Q\\
  \;&=\; (1 - c_{\eff})\,p^\top\pi \;+\; c_{\eff}\,Q,
  \qquad c_{\eff} := \tfrac{\tau v_k}{\tau v_k + \omega}.
\end{align*}

Hence $Q - p^\top\muBL = (1-c_{\eff})(Q - p^\top\pi)$. Inserting this
into $(\star)$ and simplifying the prefactor:
\begin{align*}
  \muBL \;&=\; \pi \;+\; \tfrac{\tau\,(1 - c_{\eff})}{\omega}\,(Q - p^\top\pi)\,\Sigma p\\
  \;&=\; \pi \;+\; \tfrac{\tau}{\omega + \tau v_k}\,(Q - p^\top\pi)\,\Sigma p\\
  \;&=\; \pi \;+\; c_{\eff}\,(Q - p^\top\pi)\,\tfrac{\Sigma p}{v_k}.
\end{align*}

The exact identity is established:
\begin{equation}
\label{eq:bl-single-exact}
  \muBL \;=\; \pi \;+\; c_{\eff}\,(Q - p^\top\pi)\,\frac{\Sigma p}{v_k},
  \qquad c_{\eff} \;=\; \frac{\tau v_k}{\tau v_k + \omega}.
\end{equation}
The posterior moves exactly along $\Sigma p / v_k$, normalised so that
$p^\top(\Sigma p / v_k) = 1$. Substituting \eqref{eq:bl-single-exact}
into $w_{\mathrm{BL}} = \tfrac{1}{\delta}\Sigma^{-1}\muBL$ simplifies
the weight expression as well,
\begin{equation}
\label{eq:bl-single-w}
  w_{\mathrm{BL}} \;=\; w_{\mathrm{eq}} \;+\; c_{\eff}\,(Q - p^\top\pi)\,\frac{p}{\delta\,v_k},
  \qquad w_{\mathrm{eq}} = \tfrac{1}{\delta}\Sigma^{-1}\pi.
\end{equation}
The view-driven part of $w_{\mathrm{BL}}$ lies exactly in direction
$p$, without any assumption on the volatilities.

\paragraph{Step 1: $\Delta\mu$ from $\Delta c_{\eff}$.}
Differencing \eqref{eq:bl-single-exact} at fixed $\pi$, $\Sigma$ and
$\omega$:
\begin{align*}
  \Delta\muBL
  \;&=\; \Delta c_{\eff}\,(Q - p^\top\pi)\,\tfrac{\Sigma p}{v_k}\\
  p^\top \Delta\muBL
  \;&=\; \Delta c_{\eff}\,(Q - p^\top\pi).
\end{align*}
For $p = e_i - e_j$ with $\sigma_i \approx \sigma_j$ one has
$\Sigma p \approx (v_k/p^\top p)\,p$, hence
\begin{equation}
\label{eq:dmu-vec}
  \Delta\muBL \;\approx\; \frac{\Delta c_{\eff}\,(Q - p^\top\pi)}{p^\top p}\,p,
  \qquad
  \Delta\mu_i \;\approx\; -\Delta\mu_j \;\approx\; \tfrac{1}{2}\,\Delta c_{\eff}\,(Q - p^\top\pi).
\end{equation}
This symmetric per-component split rests entirely on
$\sigma_i \approx \sigma_j$. If the volatilities differ, the direction
$\Sigma p / v_k$ distributes the return shift asymmetrically across the
two assets and the clean antisymmetric split no longer holds; the exact
vector identity \eqref{eq:bl-single-exact} is unaffected, only its
symmetric reading depends on the assumption. The numerical example below
uses a pair with nearly identical volatilities and confirms that the
symmetric split is accurate in the homogeneous case.

\paragraph{Step 2: $\Delta w$ from $\Delta\mu$.}
Differentiating \eqref{eq:bl-single-w} with respect to $c_{\eff}$ at
fixed $\Sigma$ yields directly
\begin{equation}
\label{eq:dw-dc}
  \Delta w_{\mathrm{BL}} \;=\; \frac{\Delta c_{\eff}\,(Q - p^\top\pi)}{\delta\,v_k}\,p.
\end{equation}
The shift sits exactly along $p$ with per-component
$\Delta w_i = -\Delta w_j = \Delta c_{\eff}\,(Q-p^\top\pi)/(\delta v_k)$.
\eqref{eq:dw-dc} is the exact partial derivative at fixed $\Sigma$,
i.e.\ the pure view-absorption channel. This is the quantity measured
numerically in the $2\times 2$ example below when $\hat\Sigma_s^{-1}$ is
used for both scenarios. The separate effect of a changed $\Sigma^{-1}$
is the precision channel from Section~\ref{sec:precision}.

Every percentage point of $\Delta c_{\eff}$ moves the spread
$w_i - w_j$ linearly, with leverage $1/(\delta v_k)$. Small $v_k$
amplifies the effect, and small $v_k$ is exactly the regime in which
$\Delta c_{\eff}$ is large. Affected pairs therefore combine the largest
confidence drift with the strongest translation factor into the weights.

\subsection{Real Example: Clean $2\times 2$ System with IWF and IWD}
\label{sec:example}

We isolate the effect on the expected-return vector by reducing the
universe to two assets, IWF (US large-cap growth) and IWD (US large-cap
value), and running the BL formula in closed form. No other assets
interact, so every number below is fully traceable. All quantities are
annualised and computed on the training window January 2008 to December
2018. With $\sigma_{\mathrm{IWF}} = 15.5\,\%$ and
$\sigma_{\mathrm{IWD}} = 15.6\,\%$ ($k = \sigma_i/\sigma_j = 0.99$) the
pair sits squarely in the homogeneous special case, so the symmetric
approximations of Section~\ref{sec:dmu-dw} should be accurate here.

To keep the example self-contained, the extracted $2\times 2$ sample
block is shrunk with its own target
$\mu_s^{(2)} = \tfrac{1}{2}\mathrm{tr}(\hat\Sigma_s^{(2)})$, while the
intensity $\alpha = 0.041$ is the analytic value estimated on the full
eight-asset universe. (If instead the full $8\times 8$ matrix is shrunk
and the pair extracted afterwards---the workflow behind
Table~\ref{tab:pairs}---the smaller universe-wide target
$\mu_s < \mu_s^{(2)}$ widens the spread variance further and the shift
grows to $+10.6$ percentage points. The self-contained design is thus
the conservative one.)

The inputs are
\begin{align*}
    \delta = 2.44, \qquad \tau = 0.05, \qquad c = 0.50, \qquad
    w_{\mathrm{eq}} = (0.50,\,0.50)^\top, \qquad Q = 0.03.
\end{align*}
The two annualised covariance matrices read
\begin{align*}
    \hat{\Sigma}_s &= \begin{pmatrix} 0.02408 & 0.02227 \\ 0.02227 & 0.02435 \end{pmatrix},
    \qquad \rho_s = 0.920,\\
    \hat{\Sigma}_{\LW} &= \begin{pmatrix} 0.02409 & 0.02137 \\ 0.02137 & 0.02435 \end{pmatrix},
    \qquad \rho_{\LW} = 0.882, \qquad \alpha = 0.041.
\end{align*}
$\omega$ is calibrated once on the sample estimator and then frozen,
$\omega = \tfrac{1-c}{c}\,\tau\,v_k^\Std = 1.95 \cdot 10^{-4}$.

The effective view absorption follows directly,
\begin{align*}
    v_k^\Std &= 0.00389, & v_k^\LW &= 0.00570,\\
    c_{\eff}(\hat{\Sigma}_s) &= 50.00\,\%, & c_{\eff}(\hat{\Sigma}_{\LW}) &= 59.42\,\%,\\
    & & \Delta c_{\eff} &= +9.42\,\text{pp}.
\end{align*}

The expected-return vector and the corresponding portfolio weights shift
accordingly. The LW weights deliberately use $\hat{\Sigma}_s^{-1}$ as
well, so only the view-absorption channel shows up; the separate
precision channel through $\hat{\Sigma}^{-1}$ is treated in
Section~\ref{sec:precision}.
\begin{align*}
    \mu_{\mathrm{BL}}^\Std &= (\,6.37\,\%,\ 4.89\,\%\,), & \mu_{\mathrm{BL}}^\LW &= (\,6.52\,\%,\ 4.75\,\%\,),\\
    w_{\mathrm{BL}}^\Std &= (\,+209.4\,\%,\ -109.4\,\%\,), & w_{\mathrm{BL}}^\LW &= (\,+239.6\,\%,\ -139.2\,\%\,).
\end{align*}
The difference vectors read
\[
    \boxed{\;
    \Delta c_{\eff} = +9.42\,\text{pp}, \qquad
    \Delta\mu = (+0.15\,\text{pp},\ -0.13\,\text{pp}), \qquad
    \Delta w = (+30.2\,\text{pp},\ -29.8\,\text{pp}).
    \;}
\]

The closed-form formulas from the previous subsection reproduce these
numbers. With $p^\top p = 2$, $Q - p^\top\pi = 0.0303$, $\delta = 2.44$
and $v_k = 0.00389$ from the sample estimator we check
\begin{align*}
    \Delta\mu_{\mathrm{IWF}} \;&\approx\; \tfrac{1}{2}\,\Delta c_{\eff}\,(Q-p^\top\pi)
    \;=\; \tfrac{1}{2}\cdot 0.0942 \cdot 0.0303
    \;=\; +0.143\,\text{pp},\\
    \Delta w_{\mathrm{IWF}} \;&\approx\; \frac{\Delta c_{\eff}\,(Q-p^\top\pi)}{\delta\,v_k}
    \;=\; \frac{0.0942 \cdot 0.0303}{2.44 \cdot 0.00389}
    \;=\; +0.301
    \;=\; +30.1\,\text{pp}.
\end{align*}
The IWD counterparts have the opposite sign. Both predictions match the
numerical $\Delta\mu$ and $\Delta w$ to within a few hundredths of a
percentage point. The remaining gap of about $0.01\,\text{pp}$ in
$\Delta\mu$ stems from the slight asymmetry
$\sigma_{\mathrm{IWF}} \neq \sigma_{\mathrm{IWD}}$ that the
approximation $\Sigma p / v_k \approx p / (p^\top p)$ ignores.

The chain runs in one direction. The correlation between IWF and IWD
falls from $0.920$ to $0.882$ under shrinkage, $v_k$ grows from
$0.00389$ to $0.00570$, and $c_{\eff}$ climbs by $9.42$ percentage
points. The posterior pushes IWF up by $0.15$ pp and IWD down by $0.13$
pp. Through $w = \tfrac{1}{\delta}\Sigma^{-1}\mu$ this return shift
translates into roughly $\pm 30$ percentage points on the long-short
spread. The analyst still has $c = 50\,\%$ on the input form. The model
is operating at $c_{\eff} = 59.42\,\%$.

\subsection{Implementation in Pseudocode}
\label{sec:pseudocode}

The whole computation reduces to a few lines. Inputs are the sample
covariance, the analytic Ledoit--Wolf intensity, the market weights,
$\delta$, $\tau$, and a single relative view $(p, Q, c)$. The function
\texttt{mu\_bl} solves the linear system from the Black--Litterman
first-order condition; weights are computed with $\hat\Sigma_s^{-1}$ in
both scenarios to isolate the view-absorption channel.

\begin{lstlisting}
# ----- Inputs -----
Sigma_s, alpha, w_eq, delta, tau    # market and estimator parameters
p, Q, c                             # single relative view: p = e_i - e_j

# ----- Ledoit-Wolf estimator -----
n        = dim(Sigma_s)
mu_s     = trace(Sigma_s) / n
Sigma_lw = (1 - alpha) * Sigma_s + alpha * mu_s * I_n

# ----- View variance under both estimators -----
v_std = p.T @ Sigma_s  @ p
v_lw  = p.T @ Sigma_lw @ p

# ----- Idzorek omega, frozen on Sample-Sigma -----
omega = ((1 - c) / c) * tau * v_std

# ----- Effective view absorption -----
c_eff_std = tau * v_std / (tau * v_std + omega)   # equals c by construction
c_eff_lw  = tau * v_lw  / (tau * v_lw  + omega)

# ----- BL posterior (single relative view) -----
pi = delta * Sigma_s @ w_eq

def mu_bl(Sigma):
    M = inv(tau * Sigma) + outer(p, p) / omega
    b = inv(tau * Sigma) @ pi + p * Q / omega
    return solve(M, b)

mu_BL_std = mu_bl(Sigma_s)
mu_BL_lw  = mu_bl(Sigma_lw)

# ----- Weights (return-channel isolation: same Sigma^-1) -----
w_BL_std = solve(Sigma_s, mu_BL_std) / delta
w_BL_lw  = solve(Sigma_s, mu_BL_lw)  / delta

# ----- Output -----
delta_c_eff = c_eff_lw - c_eff_std      # scalar
delta_mu    = mu_BL_lw - mu_BL_std      # (n,) vector
delta_w     = w_BL_lw  - w_BL_std       # (n,) vector
\end{lstlisting}

\noindent
On the IWF/IWD inputs from Section~\ref{sec:example} the routine returns
\texttt{delta\_c\_eff = +0.0942}, \texttt{delta\_mu = (+0.0015, -0.0013)}
and \texttt{delta\_w = (+0.302, -0.298)}, matching the boxed result.

\subsection{Recommendation}
\label{sec:rec}

The clean fix is to recompute $\omega_k$ alongside $\Sigma$,
\[
  \omega_k^{\LW} \;=\; \frac{1-c_k}{c_k}\,\tau\,v_k^{\LW}.
\]
This restores $c_{\eff,k} = c_k$ exactly. Covariance estimator,
confidence and view-uncertainty form one unit. Replacing one of the
three pieces without re-deriving the other two is the source of the
leakage, not the choice of estimator itself. The point is not academic:
the analyst who freezes $\Omega$ has, without knowing it, raised the
confidence of every affected view---a breach of whatever risk-governance
process produced the original confidences.

% ============================================================
\section{Conditioning and Precision Channel}
\label{sec:precision}
% ============================================================

The Idzorek case showed that the view absorption stays stable. The
fixed-$\Omega$ case showed that the shift in the return signal becomes
material under identifiable conditions. The larger effect of the LW
switch, however, lies not in the return estimator but in the inverse
$\hat\Sigma^{-1}$ entering
$w_{\mathrm{BL}} = \frac{1}{\delta}\Sigma^{-1}\mu_{\mathrm{BL}}$. The
weight change decomposes into two channels,
\[
    \Delta w_{\mathrm{BL}}
    \;=\; \underbrace{\frac{1}{\delta}\hat{\Sigma}_{\LW}^{-1}
          \bigl(\mu_{\mathrm{BL}}^{\LW} - \mu_{\mathrm{BL}}^{\Std}\bigr)
         }_{\text{return channel}}
    \;+\; \underbrace{\frac{1}{\delta}
          \bigl(\hat{\Sigma}_{\LW}^{-1} - \hat{\Sigma}_s^{-1}\bigr)\,
          \mu_{\mathrm{BL}}^{\Std}
         }_{\text{precision channel}}.
\]
The return channel is small because Idzorek fixes view absorption at $c$
and $\Delta\pi = O(\alpha)$ is the only remaining term. The precision
channel is driven by the difference of inverses.

LW lifts every eigenvalue $\lambda_k$ of $\hat\Sigma_s$,
\[
    \lambda_k^{\LW} \;=\; (1-\alpha)\,\lambda_k \;+\; \alpha\,\mu_s,
    \qquad k = 1,\dots,n.
\]
The smallest eigenvalue rises by the largest factor, so
$1/\lambda_n^{\LW} \ll 1/\lambda_n$. As a result $\hat\Sigma_{\LW}^{-1}$
is much more uniform than $\hat\Sigma_s^{-1}$. Extreme long-short
positions, which the badly conditioned estimator produces along
low-variance directions, get pulled back. On the Idzorek universe the
condition number drops from $341.8$ to $108.0$; under the three-view
setup of Section~\ref{sec:backtest}, the IWF position (US large-cap
growth) falls from $+170.2\,\%$ to $+119.9\,\%$ and the IWD short
(US large-cap value) from $-133.1\,\%$ to $-81.7\,\%$, taking the gross
exposure from $4.20$ to $3.10$.

% ============================================================
\section{Out-of-Sample Backtest}
\label{sec:backtest}
% ============================================================

All weights are estimated once on the training window (January 2008 to
December 2018) and held constant on the test window (January 2019 to
December 2024). The three relative views of Section~\ref{sec:blm-brief}
apply: IWF $>$ IWD by $+3\,\%$ p.a., EEM $>$ EFA by $+4\,\%$ p.a., IWF
$>$ AGG by $+6\,\%$ p.a., each at confidence $c = 0.50$, $\tau = 0.05$,
with Idzorek-calibrated $\Omega$ (recomputed per estimator, so the view
channel is clean by construction). The set follows Idzorek's example
views in structure---relative pair views with moderate spreads---and is
deliberately not tuned to the test window; ex post, two views turn out
right and one wrong. All performance figures are computed on monthly
excess returns.

\begin{table}[H]
\centering
\renewcommand{\arraystretch}{1.15}
\caption{Out-of-sample performance (Jan.\,2019--Dec.\,2024), Idzorek
universe, buy-and-hold weights.
$\kappa(\hat\Sigma_s) = 341.8$, $\kappa(\hat\Sigma_{\LW}) = 108.0$.}
\label{tab:backtest}
\small
\begin{tabular}{lrrr}
\toprule
Metric & Market equilibrium & BLM (sample) & BLM (Ledoit--Wolf) \\
\midrule
Annual return     & $3.85\,\%$   & $17.78\,\%$  & $13.46\,\%$  \\
Annual volatility & $12.16\,\%$  & $24.39\,\%$  & $19.41\,\%$  \\
Sharpe ratio      & $0.32$       & $\mathbf{0.73}$ & $0.69$    \\
Max drawdown      & $-24.78\,\%$ & $-51.30\,\%$ & $-42.26\,\%$ \\
\bottomrule
\end{tabular}
\end{table}

The largest change in any posterior mean across the two estimators is
$0.44\,\%$ p.a., of which the pure prior channel contributes at most
$0.25\,\%$ (Section~\ref{sec:idzorek-cancel}); the remainder enters
through the $\Sigma$-dependent view projection at unchanged absorption
$c$. The Idzorek cancellation works as expected, and the performance gap
is generated almost entirely by the precision channel.
$\hat\Sigma_s^{-1}$ produces $+170.2\,\%$ in US large-cap growth and
$-133.1\,\%$ in US large-cap value; $\hat\Sigma_{\LW}^{-1}$ trims these
to $+119.9\,\%$ and $-81.7\,\%$. Because growth stocks strongly
outperformed in the test window, the concentrated sample-covariance bet
pays off ($17.78\,\%$, Sharpe $0.73$). Ledoit--Wolf cuts volatility to
$19.41\,\%$ and the maximum drawdown to $-42.26\,\%$, at the cost of a
portion of the upside. Fewer extreme positions mean less risk and less
return in an up-market. With Idzorek $\Omega$ the LW switch does not
change the return signal, it only dampens the leverage of the
optimisation step.

LW underperforms the sample-covariance BL on annualised return and
Sharpe in this specific test window, a strong US growth market. The
mechanism is structural: LW caps the concentrated growth position, and
because US large-cap growth dominated every other asset class in the
test period, LW gives up return. The maximum drawdown, however, falls
from $-51.30\,\%$ to $-42.26\,\%$, a nine-point reduction. This is the
dimension where LW carries its own weight; a drawdown above $50\,\%$ is
simply not investable for many institutional mandates, regardless of the
year-end print. Imposing a long-only constraint instead compresses both
variants toward each other and toward far tamer risk figures, consistent
with the observation of \citet[p.~1652]{JagannathanMa2003} that
no-short-sale constraints act on the same estimation-error leverage that
shrinkage targets; investors who optimise long-only to begin with gain
correspondingly little from the estimator switch.

When is LW preferable? Three configurations matter. Regulated portfolios
with explicit drawdown or VaR constraints benefit, because the loss
limit, not the risk-adjusted return, is the binding constraint. Periods
of high market volatility or correlation regime shifts are a second
case, because badly conditioned sample estimators tend to produce
spurious extreme positions, especially after short training windows or
at low $T/n$. And third, whenever the training window contains a
structural regime anomaly (such as the unusually wide eigenvalue spreads
of 2008--2009) that flows into the sample covariance and propagates into
the test window.

% ============================================================
\section{Conclusion}
\label{sec:conclusion}
% ============================================================

Whether the Ledoit--Wolf switch is harmless inside Black--Litterman is
decided by the treatment of $\Omega$, and both regimes now have closed
forms. Under Idzorek's calibration
$\omega_k(\Sigma) = \tfrac{1-c_k}{c_k}\tau\,p_k^\top\Sigma p_k$ the
absorption $c_{\eff,k}$ is independent of $\Sigma$: any positive-definite
estimator can be plugged in without distorting the view channel, and the
posterior impact is confined to the $O(\alpha)$ prior shift. If
$\omega_k$ is instead computed once and frozen, the estimator switch
shifts the effective confidence by the exact, $\tau$-free expression
\eqref{eq:dceff-exact}. The sign criterion is
$v_k^\Std \lessgtr 2\mu_s$, the symmetric special case yields the
closed-form correlation threshold \eqref{eq:rho-star}, and the exact
single-view identity \eqref{eq:bl-single-exact} propagates the shift
into posterior means and weights with leverage $1/(\delta v_k)$. On the
eight-asset Idzorek universe at $\alpha = 0.041$, five of 28 pairs cross
the five-percentage-point band, led by IWF/IWD at $+10.6$; the bond pair
AGG/BWX crosses it at a correlation of only $0.57$. Recomputing
$\omega_k$ alongside $\Sigma$ removes the leakage analytically. The
lesson is procedural: covariance, confidence and view uncertainty have
to be kept consistent with each other.

Several limitations point to extensions. The closed-form chain covers a
single relative view; with several simultaneous views the shifts
interact through the posterior precision, and the multi-view analogue
(via Sherman--Morrison updates) is work in progress. The empirical part
rests on one training/test split of a small, well-diversified ETF
universe on which the analytic intensity is modest; equity universes
with larger $n/T$ push $\alpha$, and with it the leakage, substantially
higher. Finally, the mechanism is target-independent---replacing the
scaled identity by the single-factor target of \citet{LedoitWolf2003} or
the nonlinear shrinkage of \citet{LedoitWolf2017} only changes the
expression for $v_k^{\LW} - v_k^\Std$---but the quantitative reach of
the effect under these estimators remains to be mapped out.

% ============================================================
\section*{Reproducibility}
% ============================================================

All numbers quoted in this paper are produced by the script
\texttt{code/paper\_numbers.py} from frozen monthly return data
(\texttt{data/}, sourced from Yahoo Finance), using the same estimation
pipeline as the underlying thesis code: excess returns, sample
covariance with $T-1$ normalisation, the analytic Ledoit--Wolf estimator
of \texttt{scikit-learn}, and implied risk aversion from S\&P~500 excess
returns. Code and data are available at
\url{https://github.com/lucasposern/bsc-thesis-capm-blm}.
% TODO: replace by / add Zenodo DOI 10.5281/zenodo.20787490 once v2.0 is published.

% ============================================================
\bibliographystyle{plainnat}
% TODO (journal stage): switch to BibTeX (../Quellen/quellen.bib) and the
% target journal's citation style.
\begin{thebibliography}{19}

\bibitem[Best and Grauer(1991)]{BestGrauer1991}
M.~J. Best and R.~R. Grauer.
\newblock On the sensitivity of mean-variance-efficient portfolios to
changes in asset means: Some analytical and computational results.
\newblock \emph{The Review of Financial Studies}, 4(2):315--342, 1991.

\bibitem[Black and Litterman(1992)]{BlackLitterman1992}
F.~Black and R.~Litterman.
\newblock Global portfolio optimization.
\newblock \emph{Financial Analysts Journal}, 48(5):28--43, 1992.

\bibitem[DeMiguel et~al.(2009)]{DeMiguelGarlappiUppal2009}
V.~DeMiguel, L.~Garlappi, and R.~Uppal.
\newblock Optimal versus naive diversification: How inefficient is the
$1/N$ portfolio strategy?
\newblock \emph{The Review of Financial Studies}, 22(5):1915--1953, 2009.

\bibitem[He and Litterman(1999)]{HeLitterman1999}
G.~He and R.~Litterman.
\newblock The intuition behind Black--Litterman model portfolios.
\newblock Investment Management Research, Goldman Sachs \& Co., 1999.

\bibitem[Idzorek(2007)]{Idzorek2007}
T.~Idzorek.
\newblock A step-by-step guide to the Black--Litterman model:
Incorporating user-specified confidence levels.
\newblock In S.~Satchell, editor, \emph{Forecasting Expected Returns in
the Financial Markets}, pages 17--38. Academic Press, 2007.

\bibitem[Jagannathan and Ma(2003)]{JagannathanMa2003}
R.~Jagannathan and T.~Ma.
\newblock Risk reduction in large portfolios: Why imposing the wrong
constraints helps.
\newblock \emph{The Journal of Finance}, 58(4):1651--1683, 2003.

\bibitem[Ledoit and Wolf(2003)]{LedoitWolf2003}
O.~Ledoit and M.~Wolf.
\newblock Improved estimation of the covariance matrix of stock returns
with an application to portfolio selection.
\newblock \emph{Journal of Empirical Finance}, 10(5):603--621, 2003.

\bibitem[Ledoit and Wolf(2004{\natexlab{a}})]{LedoitWolf2004}
O.~Ledoit and M.~Wolf.
\newblock A well-conditioned estimator for large-dimensional covariance
matrices.
\newblock \emph{Journal of Multivariate Analysis}, 88(2):365--411,
2004{\natexlab{a}}.

\bibitem[Ledoit and Wolf(2004{\natexlab{b}})]{LedoitWolfHoney2004}
O.~Ledoit and M.~Wolf.
\newblock Honey, I shrunk the sample covariance matrix.
\newblock \emph{The Journal of Portfolio Management}, 30(4):110--119,
2004{\natexlab{b}}.

\bibitem[Ledoit and Wolf(2017)]{LedoitWolf2017}
O.~Ledoit and M.~Wolf.
\newblock Nonlinear shrinkage of the covariance matrix for portfolio
selection: Markowitz meets Goldilocks.
\newblock \emph{The Review of Financial Studies}, 30(12):4349--4388, 2017.

\bibitem[Michaud(1989)]{Michaud1989}
R.~O. Michaud.
\newblock The Markowitz optimization enigma: Is `optimized' optimal?
\newblock \emph{Financial Analysts Journal}, 45(1):31--42, 1989.

\bibitem[Satchell and Scowcroft(2000)]{SatchellScowcroft2000}
S.~Satchell and A.~Scowcroft.
\newblock A demystification of the Black--Litterman model: Managing
quantitative and traditional portfolio construction.
\newblock \emph{Journal of Asset Management}, 1(2):138--150, 2000.

\bibitem[Walters(2014)]{Walters2014}
J.~Walters.
\newblock The Black--Litterman model in detail.
\newblock Working paper, SSRN 1314585, 2014.

\end{thebibliography}

\end{document}
